\documentclass[11pt, a4paper]{article}
\usepackage[utf8]{inputenc}
\usepackage[spanish]{babel}
\usepackage[margin=1.5cm]{geometry}
\usepackage{tabularx}
\usepackage{amssymb}
\usepackage[most]{tcolorbox}
\usepackage{booktabs}
\usepackage{enumitem}

\begin{document}
\begin{center}
    {\Large \textbf{Hoja de Calificación del Instructor - Defensa Hito 1}}\\[0.2cm]
    {\large Circuitos Electrónicos III (IMT-221)}
\end{center}
\hrule
\vspace{0.3cm}

\begin{tcolorbox}[colback=white, colframe=black, title=\textbf{1. Identificación}]
\begin{tabularx}{\textwidth}{lXlX}
\textbf{Estudiante:} & \hrulefill & \textbf{Grupo:} & \hrulefill \\
\textbf{Fecha:} & \hrulefill & \textbf{Topología Asignada:} & \hrulefill \\
\multicolumn{4}{l}{\textbf{Sustituciones de Componentes:} \hrulefill} \\
\end{tabularx}
\end{tcolorbox}

\begin{tcolorbox}[colback=white, colframe=black, title=\textbf{2. Disponibilidad de Evidencia}]
\begin{tabularx}{\textwidth}{XXX}
$\square$ Cuaderno/Python & $\square$ Hardware Funcional & $\square$ Comparación Teoría-Sim-Exp \\
$\square$ Archivo LTspice & $\square$ Osciloscopio/CSV & $\square$ Reporte/Artículo \\
\end{tabularx}
\end{tcolorbox}

\begin{tcolorbox}[colback=white, colframe=black, title=\textbf{3. Dimensiones de la Rúbrica (Puntaje Oficial)}]
\begin{tabularx}{\textwidth}{X X}
Modelado Matemático + Python (20\%): \_\_\_\_ / 20 & Simulación + Hardware (30\%): \_\_\_\_ / 30 \\
Validación + Análisis (20\%): \_\_\_\_ / 20 & Defensa Dinámica Individual (30\%): \_\_\_\_ / 30 \\
\end{tabularx}
\end{tcolorbox}

\begin{tcolorbox}[colback=blue!5, colframe=blue!75!black, title=\textbf{4. Prueba de Parámetro en Vivo (Live Prompt)}]
\textbf{Parámetro Modificado:} \hrulefill \\[0.2cm]
\textbf{Predicción del Estudiante:} \hrulefill \\[0.2cm]
\textbf{Calidad de la Predicción:} $\square$ Correcta \hspace{0.2cm} $\square$ Parcialmente correcta \hspace{0.2cm} $\square$ Sin sustento \hspace{0.2cm} $\square$ Sin predicción \\[0.2cm]
\textbf{Explicación Física:} \hrulefill \\[0.2cm]
\textbf{Ubicación de Código / Ecuación Explicada:} \hrulefill 
\end{tcolorbox}

\begin{tcolorbox}[colback=red!5, colframe=red!75!black, title=\textbf{5. Diagnóstico Obligatorio: Transitorios y Flyback}]
$\square$ Explica correctamente la continuidad de corriente ($i_L(0^-) = i_L(0^+)$)\\
$\square$ Explica correctamente la polaridad del voltaje y el mecanismo de desconexión ($v_L = L\,di/dt$)\\
$\square$ La trayectoria de corriente de recirculación Flyback es correcta\\
$\square$ La orientación del diodo está justificada físicamente
\end{tcolorbox}

\begin{tcolorbox}[colback=white, colframe=black, title=\textbf{6. Diagnóstico Final}]
\textbf{Debilidades o Conceptos Erróneos Observados:} \hrulefill \\[0.2cm]
\hrulefill \\[0.2cm]
\textbf{Estado de Recuperación:}\\
$\square$ Ninguno / Listo para BJT \hspace{0.2cm} $\square$ Refuerzo delimitado \hspace{0.2cm} $\square$ Requiere evidencia de recuperación
\end{tcolorbox}

\begin{center}
    \Large \textbf{Puntaje Final del Hito: \_\_\_\_ / 100}
\end{center}

\end{document}
